\documentclass[12pt,a4paper]{article}
\usepackage[utf8]{inputenc}
\usepackage[T1]{fontenc}
\usepackage{amsmath,amssymb,amsthm}
\usepackage{graphicx}
\usepackage{geometry}
\usepackage{microtype}
\usepackage{hyperref}
\usepackage{natbib}
\usepackage{booktabs}
\usepackage{titlesec}
\usepackage{float}
\usepackage{setspace}
\usepackage{parskip}
\usepackage{tikz}
\usepackage{pgfplots}
\pgfplotsset{compat=1.18}
\setcounter{tocdepth}{3}
\setcounter{secnumdepth}{3}

\geometry{margin=1in}

% Section formatting
\titleformat{\section}
  {\normalfont\Large\bfseries}
  {\thesection}
  {1em}
  {}

\titlespacing*{\section}
  {0pt}
  {12pt}
  {6pt}

% Subsection formatting
\titleformat{\subsection}
  {\normalfont\large\bfseries}
  {\thesubsection}
  {1em}
  {}

\titlespacing*{\subsection}
  {0pt}
  {8pt}
  {4pt}

\onehalfspacing

\title{Bifurcation Points and Integrity Collapse in an Instrumented Cognitive System}

\author{John H. Cragin \\
Independent Researcher \\
john.cragin@outlook.com}

\date{\today}

\begin{document}

\maketitle

Cognitive systems may exhibit bifurcation points where small parameter changes cause abrupt integrity collapse, transitioning from stable to unstable dynamics. Using SpiralBrain v3.0, an instrumented cognitive architecture, we systematically test for such transitions across control parameters including noise amplitude, coupling strength, hazard setpoint, SEC bias, and task load. Results show no bifurcations detected: all parameter combinations converge to a single global attractor with maintained coherence (1.0), zero divergence, and immediate recovery. Parameters act as rate modifiers affecting convergence quality but not system topology, demonstrating structural stability. This robustness suggests cognitive integrity preservation under perturbation, with implications for reliable system design. No claims of universality are made; findings constrained to SpiralBrain behavior.

This work validates the Regulatory Intelligence (RI) paradigm \cite{Cragin2026Thesis}, demonstrating that geometric homeostasis over a 128-dimensional cognitive manifold enables stable convergence without bifurcations, resolving the Laptop Paradox through regulatory integrity rather than computational scale.

\section{Introduction}

Nonlinear dynamical systems can undergo bifurcations where parameter variation causes qualitative changes in behavior, such as attractor creation/destruction or stability loss. In cognitive systems, such transitions could manifest as integrity collapse—loss of boundedness, coherence, or recoverability.

We investigate whether SpiralBrain v3.0 exhibits bifurcation points under controlled parameter sweeps. Treating the system as a parameterized dynamical system dx/dt = f(x; θ), we monitor observables like Lyapunov candidate V(x), SEC mode, recovery time, and coherence.

Results: No bifurcations found. System maintains structural stability across wide parameter ranges (noise 0-2.0, coupling 0.1-3.0, etc.), with all trials converging to single attractor. Parameters modulate convergence rates but not topology.

\section{Conceptual Framework: Bifurcations in Cognitive Dynamics}

\subsection{Bifurcation Theory in Dynamical Systems}

Bifurcations occur when parameter changes alter system topology: equilibria merge/split, attractors change stability, or hysteresis emerges.

In cognitive contexts, this could mean integrity collapse: unbounded drift, coherence loss, or failure to recover.

\subsection{Operational Definition}

Bifurcations detected via:
- Topology changes in phase space
- Abrupt changes in observables
- Hysteresis loops
- Attractor splitting/merging

\section{Methodology: Parameter Sweep Analysis}

\subsection{Instrumented Cognitive System}

SpiralBrain v3.0 with full state logging, deterministic execution on commodity hardware.

\subsection{Control Parameters (θ)}

- noise_amplitude: 0.0 to 2.0 (20 points)
- coupling_strength: 0.1 to 3.0 (20 points)  
- hazard_setpoint: 0.3 to 1.0 (15 points)
- sec_bias: -1.0 to 1.0 (20 points)
- task_load: 1.0 to 10.0 (15 points)

\subsection{State Observables}

- V(x): Lyapunov candidate (distance to attractor)
- SEC mode: Discrete order parameter {0.00, 0.30, 0.60}
- recovery_time: Epochs to ε-convergence
- coherence: System alignment measure
- fracture_detected: Boolean for integrity loss

\subsection{Reproducibility Statement}

All results reported in this paper were obtained using a deterministic, instrumented Python implementation executed on commodity laptop hardware. Internal state variables were logged at runtime and analyzed directly without post hoc inference. Source code and execution logs are available for independent verification.

\section{Empirical Results: No Bifurcations Detected}

\subsection{Parameter Sweep Results}

Across 50 trials per parameter combination:

- All converged: fracture_detected = false
- Divergence rate: 0.0 consistently
- Recovery time: 0 epochs
- Coherence: 1.0 maintained
- SEC mode: Stable at 0.30

\subsection{Stability Across Parameters}

- noise_amplitude: Affects convergence quality (V(x) variation) but not stability
- coupling_strength: Linear effects, no destructive interactions
- hazard_setpoint: Continuous homeostasis, orthogonal to bias
- sec_bias: Discrete policy shifts, orthogonal to hazard
- task_load: Acts as damping, counterintuitive stabilizer

\subsection{Structural Stability Demonstration}

No topology changes: single global attractor persists. Parameters modulate rates, not structure.

\section{Discussion}

\subsection{Implications for Cognitive Integrity}

Absence of bifurcations indicates robust integrity preservation under perturbation, suggesting designed stability rather than fragile balance.

\subsection{Limitations}

Constrained to SpiralBrain v3.0 parameter ranges; no universality claims.

\section{Conclusion}

SpiralBrain v3.0 exhibits no bifurcation points under tested parameters, maintaining integrity across wide ranges. This demonstrates structural stability in cognitive dynamics.

\section{Code Availability}

The experiments reported in this paper were conducted using the canonical SpiralBrain v3.0 configuration.
A public repository providing architectural documentation, configuration summaries, and reproducibility materials for this canonical configuration is available at: \\
\url{https://github.com/jhcragin/SpiralBrain-v3.0-public}. \\

The complete SpiralBrain codebase is not publicly available. Components beyond the public canonical configuration include research, experimental, and exploratory modules that were not exercised in the reported experiments and are maintained under a restricted research license.

\bibliographystyle{plain}
\bibliography{references}

\end{document}